
% ============================
% Codificación y lenguaje
% ============================
\usepackage[utf8]{inputenc}
\usepackage[T1]{fontenc}
\usepackage[spanish]{babel}

% ============================
% Colores (cargar temprano para evitar choques)
% ============================
\usepackage[dvipsnames,table,xcdraw]{xcolor}

% ============================
% Márgenes y alineación
% ============================
\usepackage[left=.75in, right=.8in, top=1in, bottom=1in]{geometry}
\usepackage{ragged2e} % \justifying, etc.
\setlength{\headheight}{14.49998pt}
\setlength{\parindent}{0pt}
\setlength{\parskip}{0.75em}
% ============================
% Matemáticas
% ============================
\usepackage{amsmath, amssymb, amsfonts, amsthm}
\usepackage{mathrsfs}
\usepackage{soul}

\newtheoremstyle{caso_style}
  {3pt}           % Espacio vertical arriba
  {3pt}           % Espacio vertical abajo
  {\normalfont}   % Fuente del cuerpo del texto (usa \itshape si prefieres cursiva)
  {}              % Sangría (vacío = sin sangría)
  {\bfseries}     % Fuente del encabezado (Negrita)
  {.}             % Puntuación al final del encabezado (un punto)
  {0.5em}         % Espacio horizontal después del encabezado
  {\thmname{#1}\thmnote{ #3}}

\theoremstyle{caso_style}
\newtheorem*{caso}{Caso}


% ============================
% Gráficos y diagramas
% ============================
\usepackage{graphicx}
\usepackage{tikz-cd}
\usepackage{tikz}
\usetikzlibrary{automata, arrows.meta, positioning}

% ============================
% Estilo y utilidades
% ============================
\usepackage{fancyhdr}
\usepackage{booktabs}
\usepackage{float}
\usepackage{enumitem}
\usepackage{multicol}
\usepackage{tcolorbox}
\usepackage{pifont}
\usepackage{bbding}
\usepackage{comment}
% \usepackage{url} % NO es necesario si cargas hyperref (lo hace internamente)

% ============================
% Bibliografía APA con biblatex
% ============================
%\usepackage{csquotes} 
%\usepackage[backend=biber,style=ieee]{biblatex} % puedes cambiar ieee por apa, numeric, etc.
%\addbibresource{referencias.bib}

% ============================
% Algoritmos (una sola vez)
% ============================
\usepackage{algorithm}
\usepackage{algpseudocode}
\renewcommand{\thealgorithm}{} % Sin numeración (opcional)

% ============================
% Definición de colores personalizados
% ============================
\definecolor{mygreen}{rgb}{0.855, 0.969, 0.651}
\definecolor{myblue}{rgb}{0.68, 0.85, 0.9}
\definecolor{mythistle}{rgb}{0.847,0.749,0.847}
\definecolor{myorange}{rgb}{1.0, 0.6, 0.2}
\definecolor{mypurple}{rgb}{0.7, 0.2, 0.9}
\definecolor{myyellow}{rgb}{1.0, 1.0, 0.6}
\definecolor{myred}{rgb}{0.9, 0.3, 0.3}
\definecolor{mycyan}{rgb}{0.6, 1.0, 1.0}
\definecolor{mybrown}{rgb}{0.6, 0.3, 0.2}
\definecolor{mygray}{rgb}{0.85, 0.85, 0.85}
\definecolor{mymagenta}{rgb}{1.0, 0.5, 1.0}

% ============================
% tcolorbox: caja para pseudocódigo (tu estilo)
% ============================
\newtcolorbox{pseudocode}{
  colback=white, colframe=black, fontupper=\ttfamily,
  boxrule=0.5mm, arc=3mm,
  left=1mm, right=1mm, top=1mm, bottom=1mm
}

% ============================
% Listado de código (listings) con colores tipo terminal
% ============================
\definecolor{terminalbackground}{RGB}{0,0,0}
\definecolor{terminaltext}{RGB}{255,255,255}
\definecolor{terminalrule}{RGB}{0,128,0}
\definecolor{terminalnumber}{RGB}{128,128,128}

\usepackage{listings}
\lstdefinestyle{terminal}{
  backgroundcolor=\color{terminalbackground},
  basicstyle=\color{terminaltext}\large\bfseries\ttfamily,
  frame=trbl,
  framesep=2pt,
  framerule=0pt,
  xleftmargin=10pt,
  xrightmargin=10pt,
  rulecolor=\color{terminalrule},
  numberstyle=\tiny\color{terminalnumber},
  showstringspaces=false,
  upquote=true
}

% ============================
% Hipervínculos (cargar casi al final)
% ============================
\usepackage{hyperref}
\hypersetup{
  colorlinks=true,
  linkcolor=Orchid,
  filecolor=magenta,
  urlcolor=Orchid,
  citecolor=Orchid,
}

% ============================
% Figuras/tablas: títulos
% ============================
\usepackage{caption}